%% ============================================================================
%% Section: Indexing and Query Processing  (CORE)
%% Included from main.tex via %% ============================================================================
%% Section: Indexing and Query Processing  (CORE)
%% Included from main.tex via %% ============================================================================
%% Section: Indexing and Query Processing  (CORE)
%% Included from main.tex via %% ============================================================================
%% Section: Indexing and Query Processing  (CORE)
%% Included from main.tex via \input{sections/indexing}
%% Owns Table tab:complexity and Fig fig:index.
%% Design faithful to the Tier-2 layout planned in CompiledCausalGraph.java and
%% the algorithms in DefaultCausalGraphRuntime.java / PathAnalyzer.java /
%% ReachabilityIndex.java. See the status note (\NOTE) on implementation state.
%% ============================================================================
\section{Indexing and Query Processing}\label{sec:index}

This section is the core of the paper. It presents the physical layout the causal
snapshot compiles to (\S\ref{sec:index:layout}), the precomputed reachability
bitsets that accelerate every structural query (\S\ref{sec:index:bitsets}), and
the query-processing algorithms built on them: reachability and path roles
(\S\ref{sec:index:reach}), conditioning-aware d-separation
(\S\ref{sec:index:dsep}), and adjustment-set enumeration
(\S\ref{sec:index:adj}). Table~\ref{tab:complexity} summarizes the per-operator
cost against a map-of-sets baseline.

\NOTE{Implementation status (remove before CR; kept here per the honest-gaps
brief). The algorithms in \S\ref{sec:index:reach}--\S\ref{sec:index:adj} are
implemented and unit-tested (Bayes-Ball d-separation, backdoor adjustment
enumeration, path roles). The \emph{dense CSR + bitset} physical layout of
\S\ref{sec:index:layout}--\S\ref{sec:index:bitsets} is specified but not yet the
running representation: the current prototype answers these queries over a
map-of-sets adjacency with DFS closures (the $O(V{+}E)$ baseline column of
Table~\ref{tab:complexity}). The bitset layout is what the \S\ref{sec:eval}
speedups will measure; the ``Index query'' column states the design target, not a
measured result. Do not claim the speedup as achieved until the eval lands.}

\subsection{Physical layout}\label{sec:index:layout}

The snapshot's variables are assigned dense integer indices $0..n{-}1$ via a
\textsf{name}$\to$\textsf{index} map, so every downstream structure is an array
rather than a hash map (Fig.~\ref{fig:index}). Edges are stored as compressed
adjacency: a forward and a reverse \emph{CSR}-style array pair
($\textsf{forwardAdj}[i]$ lists the child indices of $i$, $\textsf{reverseAdj}[i]$
its parents), giving cache-friendly, allocation-free neighbor iteration. The
topological order is stored as an \texttt{int[]} $\textsf{topoOrder}$, and each
node carries its \emph{layer} $\textsf{topoLayer}[i]$ (0 for sources), which lets
a traversal prune any neighbor that cannot lie on a path between two nodes by a
simple layer comparison. This replaces the baseline representation---a
\lstinline[style=causaldsl]|Map<String,Set<String>>| adjacency with
\texttt{String}-keyed lookups and per-query DFS---whose every structural query
costs an $O(V{+}E)$ traversal with hashing overhead.

\subsection{Reachability bitsets}\label{sec:index:bitsets}

On top of the dense layout we precompute, for every variable, the set of its
ancestors and the set of its descendants as fixed-width bitsets
$\textsf{ancestorBits}[i],\textsf{descendantBits}[i]\in\{0,1\}^{n}$, packed into
$\lceil n/64\rceil$ machine words. The ancestor closures are filled in a single
reverse-topological sweep using the recurrence
\[
  \textsf{ancestorBits}[i] \;=\; \bigcup_{p\,\in\,\mathrm{pa}(i)}
      \bigl(\textsf{ancestorBits}[p] \;\cup\; \{p\}\bigr),
\]
and descendants symmetrically in forward-topological order; each is one
$O(V{+}E)$ pass at compile time whose cost is amortized over all later queries.
The closures occupy $O(V^2/64)$ words---the deliberate space-for-time trade that
turns reachability from traversal into a word-parallel bit test. Sparse or very
large DAGs can instead keep only the CSR layout and answer reachability by a
layer-pruned DFS; \S\ref{sec:eval} characterizes the memory/latency crossover
between the two.

\begin{figure}[t]\centering
  \resizebox{\linewidth}{!}{\input{figures/index-layout}}
  \caption{Compiled index layout for a small fragment. Variables map to dense
  indices; forward/reverse CSR arrays store adjacency; a topological order and
  per-node layer drive pruning; and per-node ancestor/descendant bitsets
  ($\lceil n/64\rceil$ words each) are filled in reverse-/forward-topological
  order. Structural queries become word-parallel bit operations over the
  bitsets.}
  \label{fig:index}
\end{figure}

\subsection{Reachability and path-role queries}\label{sec:index:reach}

With the closures in hand the reachability operators are immediate:
$\mathrm{an}(X)$ and $\mathrm{de}(X)$ are a single bitset read, and
$\textsf{hasDirectedPath}(X,Y)$ is one bit test---is $Y$ set in
$\textsf{descendantBits}[X]$---in $O(1)$. The path-role operators reduce to
bitset algebra over these closures. Confounders of $T$ and $Y$ are the common
ancestors $\textsf{ancestorBits}[T]\cap\textsf{ancestorBits}[Y]$; mediators of
$T\!\to\!Y$ are $\textsf{descendantBits}[T]\cap\textsf{ancestorBits}[Y]$; and
colliders are the nodes whose parents include one reachable from $X$ and one from
$Y$ while the node itself ancestors neither---again a handful of intersections
per candidate. Each is $O(V/64)$ word operations rather than the baseline's
repeated DFS.

\subsection{Conditioning-aware d-separation}\label{sec:index:dsep}

D-separation is \emph{not} reachability: adding a variable to the conditioning set
$Z$ can block one path and, if that variable is a collider or an ancestor of one,
unblock another. The test must therefore follow edge orientation and the role of
each node under $Z$, which we do with the Bayes-Ball procedure. A search
propagates ``balls'' from $X$, tracking for each visited node whether it was
entered from a parent or from a child; at each node the admissible continuations
depend on whether the node is in $Z$: an observed non-collider blocks a ball
arriving along the chain/fork, while a collider (or a node with an observed
descendant) passes a ball arriving from two children. $X$ and $Y$ are
d-separated given $Z$ exactly when no ball reaches $Y$. The procedure is
conditioning-aware by construction and, with conditioning and frontier sets
represented as bitsets, its per-step membership tests are $O(1)$ and the whole
test $O(V/64)$ in the frontier. A single flag---cutting the outgoing edges of the
treatment---turns the same routine into the backdoor test used next.

\subsection{Adjustment-set enumeration}\label{sec:index:adj}

Identification asks whether the effect of $T$ on $Y$ is estimable by covariate
adjustment and, if so, for a valid set. A set $Z$ satisfies the \emph{backdoor
criterion} when it contains no descendant of $T$ and blocks every backdoor path
from $T$ to $Y$; the second condition is exactly the d-separation test of
\S\ref{sec:index:dsep} with $T$'s outgoing edges cut. Enumeration first restricts
the candidate pool to \emph{eligible} variables---those that are not $T$, not $Y$,
and (via $\textsf{descendantBits}[T]$) not descendants of $T$, so the
no-descendant clause holds for free. It then searches candidate sets in
increasing size and returns the first non-empty level of blocking sets, yielding
\emph{minimal} adjustment sets; the empty set is tested first, so a query with no
open backdoor path returns immediately. This is the most traversal-intensive
operator in the catalog---it issues one backdoor (d-separation) test per
candidate set---which is precisely why casting the inner test as an $O(V/64)$
bitset operation, rather than a graph traversal, is what makes identification
scale (R2).

\begin{table}[t]\centering
\caption{Per-operator cost: map-of-sets $+$ DFS baseline vs.\ the bitset index.
$k$ is the number of candidate sets an identification query tests.}
\label{tab:complexity}
\begin{tabular}{@{}lccc@{}}
\toprule
Operator & Baseline query & Index query & Index space \\
\midrule
ancestors / descendants & $O(V{+}E)$ & $O(V/64)$ & $O(V^2/64)$ \\
hasDirectedPath         & $O(V{+}E)$ & $O(1)$     & --- \\
mediators/conf./coll.   & $O(V{+}E)$ & $O(V/64)$  & --- \\
d-separation            & $O(V{+}E)$ & $O(V/64)$  & --- \\
adjustment sets         & $O(k(V{+}E))$ & $O(k\,V/64)$ & --- \\
\bottomrule
\end{tabular}
\end{table}

%% Owns Table tab:complexity and Fig fig:index.
%% Design faithful to the Tier-2 layout planned in CompiledCausalGraph.java and
%% the algorithms in DefaultCausalGraphRuntime.java / PathAnalyzer.java /
%% ReachabilityIndex.java. See the status note (\NOTE) on implementation state.
%% ============================================================================
\section{Indexing and Query Processing}\label{sec:index}

This section is the core of the paper. It presents the physical layout the causal
snapshot compiles to (\S\ref{sec:index:layout}), the precomputed reachability
bitsets that accelerate every structural query (\S\ref{sec:index:bitsets}), and
the query-processing algorithms built on them: reachability and path roles
(\S\ref{sec:index:reach}), conditioning-aware d-separation
(\S\ref{sec:index:dsep}), and adjustment-set enumeration
(\S\ref{sec:index:adj}). Table~\ref{tab:complexity} summarizes the per-operator
cost against a map-of-sets baseline.

\NOTE{Implementation status (remove before CR; kept here per the honest-gaps
brief). The algorithms in \S\ref{sec:index:reach}--\S\ref{sec:index:adj} are
implemented and unit-tested (Bayes-Ball d-separation, backdoor adjustment
enumeration, path roles). The \emph{dense CSR + bitset} physical layout of
\S\ref{sec:index:layout}--\S\ref{sec:index:bitsets} is specified but not yet the
running representation: the current prototype answers these queries over a
map-of-sets adjacency with DFS closures (the $O(V{+}E)$ baseline column of
Table~\ref{tab:complexity}). The bitset layout is what the \S\ref{sec:eval}
speedups will measure; the ``Index query'' column states the design target, not a
measured result. Do not claim the speedup as achieved until the eval lands.}

\subsection{Physical layout}\label{sec:index:layout}

The snapshot's variables are assigned dense integer indices $0..n{-}1$ via a
\textsf{name}$\to$\textsf{index} map, so every downstream structure is an array
rather than a hash map (Fig.~\ref{fig:index}). Edges are stored as compressed
adjacency: a forward and a reverse \emph{CSR}-style array pair
($\textsf{forwardAdj}[i]$ lists the child indices of $i$, $\textsf{reverseAdj}[i]$
its parents), giving cache-friendly, allocation-free neighbor iteration. The
topological order is stored as an \texttt{int[]} $\textsf{topoOrder}$, and each
node carries its \emph{layer} $\textsf{topoLayer}[i]$ (0 for sources), which lets
a traversal prune any neighbor that cannot lie on a path between two nodes by a
simple layer comparison. This replaces the baseline representation---a
\lstinline[style=causaldsl]|Map<String,Set<String>>| adjacency with
\texttt{String}-keyed lookups and per-query DFS---whose every structural query
costs an $O(V{+}E)$ traversal with hashing overhead.

\subsection{Reachability bitsets}\label{sec:index:bitsets}

On top of the dense layout we precompute, for every variable, the set of its
ancestors and the set of its descendants as fixed-width bitsets
$\textsf{ancestorBits}[i],\textsf{descendantBits}[i]\in\{0,1\}^{n}$, packed into
$\lceil n/64\rceil$ machine words. The ancestor closures are filled in a single
reverse-topological sweep using the recurrence
\[
  \textsf{ancestorBits}[i] \;=\; \bigcup_{p\,\in\,\mathrm{pa}(i)}
      \bigl(\textsf{ancestorBits}[p] \;\cup\; \{p\}\bigr),
\]
and descendants symmetrically in forward-topological order; each is one
$O(V{+}E)$ pass at compile time whose cost is amortized over all later queries.
The closures occupy $O(V^2/64)$ words---the deliberate space-for-time trade that
turns reachability from traversal into a word-parallel bit test. Sparse or very
large DAGs can instead keep only the CSR layout and answer reachability by a
layer-pruned DFS; \S\ref{sec:eval} characterizes the memory/latency crossover
between the two.

\begin{figure}[t]\centering
  \resizebox{\linewidth}{!}{%% ============================================================================
%% Figure: compiled index layout (dense CSR + reachability bitsets)
%% Included from sections/indexing.tex inside a figure env via \resizebox{\input}.
%% Requires: \usepackage{tikz}; \usetikzlibrary{positioning,fit}
%% Example DAG: A(0)->B(1), A(0)->C(2), B(1)->D(3), C(2)->D(3)
%% ============================================================================
\begin{tikzpicture}[
    font=\scriptsize,
    >={Stealth[length=1.6mm]},
    cell/.style={draw, minimum width=5.2mm, minimum height=5mm, inner sep=1pt,
                 align=center, fill=white},
    lab/.style={anchor=east, font=\scriptsize\itshape},
    node distance=2mm,
]

% ---- tiny example DAG (top-left) ----
\node[draw, circle, minimum size=5mm] (A) at (0,0)      {A};
\node[draw, circle, minimum size=5mm] (B) at (1.1,0.7)  {B};
\node[draw, circle, minimum size=5mm] (C) at (1.1,-0.7) {C};
\node[draw, circle, minimum size=5mm] (D) at (2.2,0)    {D};
\draw[->] (A) -- (B); \draw[->] (A) -- (C);
\draw[->] (B) -- (D); \draw[->] (C) -- (D);
\node[font=\scriptsize\itshape] at (1.1,-1.6) {DAG (indices A{=}0,B{=}1,C{=}2,D{=}3)};

% ---- name -> index map ----
\node[lab] (m0) at (5.2, 1.0) {nameToIndex:};
\node[cell] (mA) at (6.1,1.0) {A:0};
\node[cell, right=of mA] (mB) {B:1};
\node[cell, right=of mB] (mC) {C:2};
\node[cell, right=of mC] (mD) {D:3};

% ---- forward CSR ----
\node[lab] (f0) at (5.2, 0.3) {forwardAdj:};
\node[cell] (f1) at (6.1,0.3) {1,2};
\node[cell, right=of f1] (f2) {3};
\node[cell, right=of f2] (f3) {3};
\node[cell, right=of f3] (f4) {--};

% ---- reverse CSR ----
\node[lab] (r0) at (5.2, -0.4) {reverseAdj:};
\node[cell] (r1) at (6.1,-0.4) {--};
\node[cell, right=of r1] (r2) {0};
\node[cell, right=of r2] (r3) {0};
\node[cell, right=of r3] (r4) {1,2};

% ---- topo order / layer ----
\node[lab] (t0) at (5.2, -1.1) {topoOrder:};
\node[cell] (t1) at (6.1,-1.1) {0};
\node[cell, right=of t1] (t2) {1};
\node[cell, right=of t2] (t3) {2};
\node[cell, right=of t3] (t4) {3};
\node[lab] (l0) at (5.2, -1.7) {topoLayer:};
\node[cell] (l1) at (6.1,-1.7) {0};
\node[cell, right=of l1] (l2) {1};
\node[cell, right=of l2] (l3) {1};
\node[cell, right=of l3] (l4) {2};

% ---- ancestor bitsets (rows = variable, columns = bit for A B C D) ----
\node[font=\scriptsize\itshape] at (1.4,-2.6) {ancestorBits[i] (bit set = ancestor)};
\node[lab] at (-0.3,-3.2) {A}; \node[lab] at (-0.3,-3.7) {B};
\node[lab] at (-0.3,-4.2) {C}; \node[lab] at (-0.3,-4.7) {D};
% rows explicitly (A,B,C,D) x (A,B,C,D)
\foreach \row/\y/\vals in {A/-3.2/{0/0,1/0,2/0,3/0}, B/-3.7/{0/1,1/0,2/0,3/0},
                           C/-4.2/{0/1,1/0,2/0,3/0}, D/-4.7/{0/1,1/1,2/1,3/0}} {
  \foreach \c/\v in \vals {
    \node[cell, minimum width=4.5mm] at ({0.1+\c*0.5},\y) {\v};
  }
}

% ---- descendant bitsets ----
\node[font=\scriptsize\itshape] at (5.2,-2.6) {descendantBits[i] (bit set = descendant)};
\node[lab] at (3.6,-3.2) {A}; \node[lab] at (3.6,-3.7) {B};
\node[lab] at (3.6,-4.2) {C}; \node[lab] at (3.6,-4.7) {D};
\foreach \row/\y/\vals in {A/-3.2/{0/0,1/1,2/1,3/1}, B/-3.7/{0/0,1/0,2/0,3/1},
                           C/-4.2/{0/0,1/0,2/0,3/1}, D/-4.7/{0/0,1/0,2/0,3/0}} {
  \foreach \c/\v in \vals {
    \node[cell, minimum width=4.5mm] at ({4.0+\c*0.5},\y) {\v};
  }
}

\end{tikzpicture}
}
  \caption{Compiled index layout for a small fragment. Variables map to dense
  indices; forward/reverse CSR arrays store adjacency; a topological order and
  per-node layer drive pruning; and per-node ancestor/descendant bitsets
  ($\lceil n/64\rceil$ words each) are filled in reverse-/forward-topological
  order. Structural queries become word-parallel bit operations over the
  bitsets.}
  \label{fig:index}
\end{figure}

\subsection{Reachability and path-role queries}\label{sec:index:reach}

With the closures in hand the reachability operators are immediate:
$\mathrm{an}(X)$ and $\mathrm{de}(X)$ are a single bitset read, and
$\textsf{hasDirectedPath}(X,Y)$ is one bit test---is $Y$ set in
$\textsf{descendantBits}[X]$---in $O(1)$. The path-role operators reduce to
bitset algebra over these closures. Confounders of $T$ and $Y$ are the common
ancestors $\textsf{ancestorBits}[T]\cap\textsf{ancestorBits}[Y]$; mediators of
$T\!\to\!Y$ are $\textsf{descendantBits}[T]\cap\textsf{ancestorBits}[Y]$; and
colliders are the nodes whose parents include one reachable from $X$ and one from
$Y$ while the node itself ancestors neither---again a handful of intersections
per candidate. Each is $O(V/64)$ word operations rather than the baseline's
repeated DFS.

\subsection{Conditioning-aware d-separation}\label{sec:index:dsep}

D-separation is \emph{not} reachability: adding a variable to the conditioning set
$Z$ can block one path and, if that variable is a collider or an ancestor of one,
unblock another. The test must therefore follow edge orientation and the role of
each node under $Z$, which we do with the Bayes-Ball procedure. A search
propagates ``balls'' from $X$, tracking for each visited node whether it was
entered from a parent or from a child; at each node the admissible continuations
depend on whether the node is in $Z$: an observed non-collider blocks a ball
arriving along the chain/fork, while a collider (or a node with an observed
descendant) passes a ball arriving from two children. $X$ and $Y$ are
d-separated given $Z$ exactly when no ball reaches $Y$. The procedure is
conditioning-aware by construction and, with conditioning and frontier sets
represented as bitsets, its per-step membership tests are $O(1)$ and the whole
test $O(V/64)$ in the frontier. A single flag---cutting the outgoing edges of the
treatment---turns the same routine into the backdoor test used next.

\subsection{Adjustment-set enumeration}\label{sec:index:adj}

Identification asks whether the effect of $T$ on $Y$ is estimable by covariate
adjustment and, if so, for a valid set. A set $Z$ satisfies the \emph{backdoor
criterion} when it contains no descendant of $T$ and blocks every backdoor path
from $T$ to $Y$; the second condition is exactly the d-separation test of
\S\ref{sec:index:dsep} with $T$'s outgoing edges cut. Enumeration first restricts
the candidate pool to \emph{eligible} variables---those that are not $T$, not $Y$,
and (via $\textsf{descendantBits}[T]$) not descendants of $T$, so the
no-descendant clause holds for free. It then searches candidate sets in
increasing size and returns the first non-empty level of blocking sets, yielding
\emph{minimal} adjustment sets; the empty set is tested first, so a query with no
open backdoor path returns immediately. This is the most traversal-intensive
operator in the catalog---it issues one backdoor (d-separation) test per
candidate set---which is precisely why casting the inner test as an $O(V/64)$
bitset operation, rather than a graph traversal, is what makes identification
scale (R2).

\begin{table}[t]\centering
\caption{Per-operator cost: map-of-sets $+$ DFS baseline vs.\ the bitset index.
$k$ is the number of candidate sets an identification query tests.}
\label{tab:complexity}
\begin{tabular}{@{}lccc@{}}
\toprule
Operator & Baseline query & Index query & Index space \\
\midrule
ancestors / descendants & $O(V{+}E)$ & $O(V/64)$ & $O(V^2/64)$ \\
hasDirectedPath         & $O(V{+}E)$ & $O(1)$     & --- \\
mediators/conf./coll.   & $O(V{+}E)$ & $O(V/64)$  & --- \\
d-separation            & $O(V{+}E)$ & $O(V/64)$  & --- \\
adjustment sets         & $O(k(V{+}E))$ & $O(k\,V/64)$ & --- \\
\bottomrule
\end{tabular}
\end{table}

%% Owns Table tab:complexity and Fig fig:index.
%% Design faithful to the Tier-2 layout planned in CompiledCausalGraph.java and
%% the algorithms in DefaultCausalGraphRuntime.java / PathAnalyzer.java /
%% ReachabilityIndex.java. See the status note (\NOTE) on implementation state.
%% ============================================================================
\section{Indexing and Query Processing}\label{sec:index}

This section is the core of the paper. It presents the physical layout the causal
snapshot compiles to (\S\ref{sec:index:layout}), the precomputed reachability
bitsets that accelerate every structural query (\S\ref{sec:index:bitsets}), and
the query-processing algorithms built on them: reachability and path roles
(\S\ref{sec:index:reach}), conditioning-aware d-separation
(\S\ref{sec:index:dsep}), and adjustment-set enumeration
(\S\ref{sec:index:adj}). Table~\ref{tab:complexity} summarizes the per-operator
cost against a map-of-sets baseline.

\NOTE{Implementation status (remove before CR; kept here per the honest-gaps
brief). The algorithms in \S\ref{sec:index:reach}--\S\ref{sec:index:adj} are
implemented and unit-tested (Bayes-Ball d-separation, backdoor adjustment
enumeration, path roles). The \emph{dense CSR + bitset} physical layout of
\S\ref{sec:index:layout}--\S\ref{sec:index:bitsets} is specified but not yet the
running representation: the current prototype answers these queries over a
map-of-sets adjacency with DFS closures (the $O(V{+}E)$ baseline column of
Table~\ref{tab:complexity}). The bitset layout is what the \S\ref{sec:eval}
speedups will measure; the ``Index query'' column states the design target, not a
measured result. Do not claim the speedup as achieved until the eval lands.}

\subsection{Physical layout}\label{sec:index:layout}

The snapshot's variables are assigned dense integer indices $0..n{-}1$ via a
\textsf{name}$\to$\textsf{index} map, so every downstream structure is an array
rather than a hash map (Fig.~\ref{fig:index}). Edges are stored as compressed
adjacency: a forward and a reverse \emph{CSR}-style array pair
($\textsf{forwardAdj}[i]$ lists the child indices of $i$, $\textsf{reverseAdj}[i]$
its parents), giving cache-friendly, allocation-free neighbor iteration. The
topological order is stored as an \texttt{int[]} $\textsf{topoOrder}$, and each
node carries its \emph{layer} $\textsf{topoLayer}[i]$ (0 for sources), which lets
a traversal prune any neighbor that cannot lie on a path between two nodes by a
simple layer comparison. This replaces the baseline representation---a
\lstinline[style=causaldsl]|Map<String,Set<String>>| adjacency with
\texttt{String}-keyed lookups and per-query DFS---whose every structural query
costs an $O(V{+}E)$ traversal with hashing overhead.

\subsection{Reachability bitsets}\label{sec:index:bitsets}

On top of the dense layout we precompute, for every variable, the set of its
ancestors and the set of its descendants as fixed-width bitsets
$\textsf{ancestorBits}[i],\textsf{descendantBits}[i]\in\{0,1\}^{n}$, packed into
$\lceil n/64\rceil$ machine words. The ancestor closures are filled in a single
reverse-topological sweep using the recurrence
\[
  \textsf{ancestorBits}[i] \;=\; \bigcup_{p\,\in\,\mathrm{pa}(i)}
      \bigl(\textsf{ancestorBits}[p] \;\cup\; \{p\}\bigr),
\]
and descendants symmetrically in forward-topological order; each is one
$O(V{+}E)$ pass at compile time whose cost is amortized over all later queries.
The closures occupy $O(V^2/64)$ words---the deliberate space-for-time trade that
turns reachability from traversal into a word-parallel bit test. Sparse or very
large DAGs can instead keep only the CSR layout and answer reachability by a
layer-pruned DFS; \S\ref{sec:eval} characterizes the memory/latency crossover
between the two.

\begin{figure}[t]\centering
  \resizebox{\linewidth}{!}{%% ============================================================================
%% Figure: compiled index layout (dense CSR + reachability bitsets)
%% Included from sections/indexing.tex inside a figure env via \resizebox{\input}.
%% Requires: \usepackage{tikz}; \usetikzlibrary{positioning,fit}
%% Example DAG: A(0)->B(1), A(0)->C(2), B(1)->D(3), C(2)->D(3)
%% ============================================================================
\begin{tikzpicture}[
    font=\scriptsize,
    >={Stealth[length=1.6mm]},
    cell/.style={draw, minimum width=5.2mm, minimum height=5mm, inner sep=1pt,
                 align=center, fill=white},
    lab/.style={anchor=east, font=\scriptsize\itshape},
    node distance=2mm,
]

% ---- tiny example DAG (top-left) ----
\node[draw, circle, minimum size=5mm] (A) at (0,0)      {A};
\node[draw, circle, minimum size=5mm] (B) at (1.1,0.7)  {B};
\node[draw, circle, minimum size=5mm] (C) at (1.1,-0.7) {C};
\node[draw, circle, minimum size=5mm] (D) at (2.2,0)    {D};
\draw[->] (A) -- (B); \draw[->] (A) -- (C);
\draw[->] (B) -- (D); \draw[->] (C) -- (D);
\node[font=\scriptsize\itshape] at (1.1,-1.6) {DAG (indices A{=}0,B{=}1,C{=}2,D{=}3)};

% ---- name -> index map ----
\node[lab] (m0) at (5.2, 1.0) {nameToIndex:};
\node[cell] (mA) at (6.1,1.0) {A:0};
\node[cell, right=of mA] (mB) {B:1};
\node[cell, right=of mB] (mC) {C:2};
\node[cell, right=of mC] (mD) {D:3};

% ---- forward CSR ----
\node[lab] (f0) at (5.2, 0.3) {forwardAdj:};
\node[cell] (f1) at (6.1,0.3) {1,2};
\node[cell, right=of f1] (f2) {3};
\node[cell, right=of f2] (f3) {3};
\node[cell, right=of f3] (f4) {--};

% ---- reverse CSR ----
\node[lab] (r0) at (5.2, -0.4) {reverseAdj:};
\node[cell] (r1) at (6.1,-0.4) {--};
\node[cell, right=of r1] (r2) {0};
\node[cell, right=of r2] (r3) {0};
\node[cell, right=of r3] (r4) {1,2};

% ---- topo order / layer ----
\node[lab] (t0) at (5.2, -1.1) {topoOrder:};
\node[cell] (t1) at (6.1,-1.1) {0};
\node[cell, right=of t1] (t2) {1};
\node[cell, right=of t2] (t3) {2};
\node[cell, right=of t3] (t4) {3};
\node[lab] (l0) at (5.2, -1.7) {topoLayer:};
\node[cell] (l1) at (6.1,-1.7) {0};
\node[cell, right=of l1] (l2) {1};
\node[cell, right=of l2] (l3) {1};
\node[cell, right=of l3] (l4) {2};

% ---- ancestor bitsets (rows = variable, columns = bit for A B C D) ----
\node[font=\scriptsize\itshape] at (1.4,-2.6) {ancestorBits[i] (bit set = ancestor)};
\node[lab] at (-0.3,-3.2) {A}; \node[lab] at (-0.3,-3.7) {B};
\node[lab] at (-0.3,-4.2) {C}; \node[lab] at (-0.3,-4.7) {D};
% rows explicitly (A,B,C,D) x (A,B,C,D)
\foreach \row/\y/\vals in {A/-3.2/{0/0,1/0,2/0,3/0}, B/-3.7/{0/1,1/0,2/0,3/0},
                           C/-4.2/{0/1,1/0,2/0,3/0}, D/-4.7/{0/1,1/1,2/1,3/0}} {
  \foreach \c/\v in \vals {
    \node[cell, minimum width=4.5mm] at ({0.1+\c*0.5},\y) {\v};
  }
}

% ---- descendant bitsets ----
\node[font=\scriptsize\itshape] at (5.2,-2.6) {descendantBits[i] (bit set = descendant)};
\node[lab] at (3.6,-3.2) {A}; \node[lab] at (3.6,-3.7) {B};
\node[lab] at (3.6,-4.2) {C}; \node[lab] at (3.6,-4.7) {D};
\foreach \row/\y/\vals in {A/-3.2/{0/0,1/1,2/1,3/1}, B/-3.7/{0/0,1/0,2/0,3/1},
                           C/-4.2/{0/0,1/0,2/0,3/1}, D/-4.7/{0/0,1/0,2/0,3/0}} {
  \foreach \c/\v in \vals {
    \node[cell, minimum width=4.5mm] at ({4.0+\c*0.5},\y) {\v};
  }
}

\end{tikzpicture}
}
  \caption{Compiled index layout for a small fragment. Variables map to dense
  indices; forward/reverse CSR arrays store adjacency; a topological order and
  per-node layer drive pruning; and per-node ancestor/descendant bitsets
  ($\lceil n/64\rceil$ words each) are filled in reverse-/forward-topological
  order. Structural queries become word-parallel bit operations over the
  bitsets.}
  \label{fig:index}
\end{figure}

\subsection{Reachability and path-role queries}\label{sec:index:reach}

With the closures in hand the reachability operators are immediate:
$\mathrm{an}(X)$ and $\mathrm{de}(X)$ are a single bitset read, and
$\textsf{hasDirectedPath}(X,Y)$ is one bit test---is $Y$ set in
$\textsf{descendantBits}[X]$---in $O(1)$. The path-role operators reduce to
bitset algebra over these closures. Confounders of $T$ and $Y$ are the common
ancestors $\textsf{ancestorBits}[T]\cap\textsf{ancestorBits}[Y]$; mediators of
$T\!\to\!Y$ are $\textsf{descendantBits}[T]\cap\textsf{ancestorBits}[Y]$; and
colliders are the nodes whose parents include one reachable from $X$ and one from
$Y$ while the node itself ancestors neither---again a handful of intersections
per candidate. Each is $O(V/64)$ word operations rather than the baseline's
repeated DFS.

\subsection{Conditioning-aware d-separation}\label{sec:index:dsep}

D-separation is \emph{not} reachability: adding a variable to the conditioning set
$Z$ can block one path and, if that variable is a collider or an ancestor of one,
unblock another. The test must therefore follow edge orientation and the role of
each node under $Z$, which we do with the Bayes-Ball procedure. A search
propagates ``balls'' from $X$, tracking for each visited node whether it was
entered from a parent or from a child; at each node the admissible continuations
depend on whether the node is in $Z$: an observed non-collider blocks a ball
arriving along the chain/fork, while a collider (or a node with an observed
descendant) passes a ball arriving from two children. $X$ and $Y$ are
d-separated given $Z$ exactly when no ball reaches $Y$. The procedure is
conditioning-aware by construction and, with conditioning and frontier sets
represented as bitsets, its per-step membership tests are $O(1)$ and the whole
test $O(V/64)$ in the frontier. A single flag---cutting the outgoing edges of the
treatment---turns the same routine into the backdoor test used next.

\subsection{Adjustment-set enumeration}\label{sec:index:adj}

Identification asks whether the effect of $T$ on $Y$ is estimable by covariate
adjustment and, if so, for a valid set. A set $Z$ satisfies the \emph{backdoor
criterion} when it contains no descendant of $T$ and blocks every backdoor path
from $T$ to $Y$; the second condition is exactly the d-separation test of
\S\ref{sec:index:dsep} with $T$'s outgoing edges cut. Enumeration first restricts
the candidate pool to \emph{eligible} variables---those that are not $T$, not $Y$,
and (via $\textsf{descendantBits}[T]$) not descendants of $T$, so the
no-descendant clause holds for free. It then searches candidate sets in
increasing size and returns the first non-empty level of blocking sets, yielding
\emph{minimal} adjustment sets; the empty set is tested first, so a query with no
open backdoor path returns immediately. This is the most traversal-intensive
operator in the catalog---it issues one backdoor (d-separation) test per
candidate set---which is precisely why casting the inner test as an $O(V/64)$
bitset operation, rather than a graph traversal, is what makes identification
scale (R2).

\begin{table}[t]\centering
\caption{Per-operator cost: map-of-sets $+$ DFS baseline vs.\ the bitset index.
$k$ is the number of candidate sets an identification query tests.}
\label{tab:complexity}
\begin{tabular}{@{}lccc@{}}
\toprule
Operator & Baseline query & Index query & Index space \\
\midrule
ancestors / descendants & $O(V{+}E)$ & $O(V/64)$ & $O(V^2/64)$ \\
hasDirectedPath         & $O(V{+}E)$ & $O(1)$     & --- \\
mediators/conf./coll.   & $O(V{+}E)$ & $O(V/64)$  & --- \\
d-separation            & $O(V{+}E)$ & $O(V/64)$  & --- \\
adjustment sets         & $O(k(V{+}E))$ & $O(k\,V/64)$ & --- \\
\bottomrule
\end{tabular}
\end{table}

%% Owns Table tab:complexity and Fig fig:index.
%% Design faithful to the Tier-2 layout planned in CompiledCausalGraph.java and
%% the algorithms in DefaultCausalGraphRuntime.java / PathAnalyzer.java /
%% ReachabilityIndex.java. See the status note (\NOTE) on implementation state.
%% ============================================================================
\section{Indexing and Query Processing}\label{sec:index}

This section is the core of the paper. It presents the physical layout the causal
snapshot compiles to (\S\ref{sec:index:layout}), the precomputed reachability
bitsets that accelerate every structural query (\S\ref{sec:index:bitsets}), and
the query-processing algorithms built on them: reachability and path roles
(\S\ref{sec:index:reach}), conditioning-aware d-separation
(\S\ref{sec:index:dsep}), and adjustment-set enumeration
(\S\ref{sec:index:adj}). Table~\ref{tab:complexity} summarizes the per-operator
cost against a map-of-sets baseline.

\NOTE{Implementation status (remove before CR; kept here per the honest-gaps
brief). The algorithms in \S\ref{sec:index:reach}--\S\ref{sec:index:adj} are
implemented and unit-tested (Bayes-Ball d-separation, backdoor adjustment
enumeration, path roles). The \emph{dense CSR + bitset} physical layout of
\S\ref{sec:index:layout}--\S\ref{sec:index:bitsets} is specified but not yet the
running representation: the current prototype answers these queries over a
map-of-sets adjacency with DFS closures (the $O(V{+}E)$ baseline column of
Table~\ref{tab:complexity}). The bitset layout is what the \S\ref{sec:eval}
speedups will measure; the ``Index query'' column states the design target, not a
measured result. Do not claim the speedup as achieved until the eval lands.}

\subsection{Physical layout}\label{sec:index:layout}

The snapshot's variables are assigned dense integer indices $0..n{-}1$ via a
\textsf{name}$\to$\textsf{index} map, so every downstream structure is an array
rather than a hash map (Fig.~\ref{fig:index}). Edges are stored as compressed
adjacency: a forward and a reverse \emph{CSR}-style array pair
($\textsf{forwardAdj}[i]$ lists the child indices of $i$, $\textsf{reverseAdj}[i]$
its parents), giving cache-friendly, allocation-free neighbor iteration. The
topological order is stored as an \texttt{int[]} $\textsf{topoOrder}$, and each
node carries its \emph{layer} $\textsf{topoLayer}[i]$ (0 for sources), which lets
a traversal prune any neighbor that cannot lie on a path between two nodes by a
simple layer comparison. This replaces the baseline representation---a
\lstinline[style=causaldsl]|Map<String,Set<String>>| adjacency with
\texttt{String}-keyed lookups and per-query DFS---whose every structural query
costs an $O(V{+}E)$ traversal with hashing overhead.

\subsection{Reachability bitsets}\label{sec:index:bitsets}

On top of the dense layout we precompute, for every variable, the set of its
ancestors and the set of its descendants as fixed-width bitsets
$\textsf{ancestorBits}[i],\textsf{descendantBits}[i]\in\{0,1\}^{n}$, packed into
$\lceil n/64\rceil$ machine words. The ancestor closures are filled in a single
reverse-topological sweep using the recurrence
\[
  \textsf{ancestorBits}[i] \;=\; \bigcup_{p\,\in\,\mathrm{pa}(i)}
      \bigl(\textsf{ancestorBits}[p] \;\cup\; \{p\}\bigr),
\]
and descendants symmetrically in forward-topological order; each is one
$O(V{+}E)$ pass at compile time whose cost is amortized over all later queries.
The closures occupy $O(V^2/64)$ words---the deliberate space-for-time trade that
turns reachability from traversal into a word-parallel bit test. Sparse or very
large DAGs can instead keep only the CSR layout and answer reachability by a
layer-pruned DFS; \S\ref{sec:eval} characterizes the memory/latency crossover
between the two.

\begin{figure}[t]\centering
  \resizebox{\linewidth}{!}{%% ============================================================================
%% Figure: compiled index layout (dense CSR + reachability bitsets)
%% Included from sections/indexing.tex inside a figure env via \resizebox{\input}.
%% Requires: \usepackage{tikz}; \usetikzlibrary{positioning,fit}
%% Example DAG: A(0)->B(1), A(0)->C(2), B(1)->D(3), C(2)->D(3)
%% ============================================================================
\begin{tikzpicture}[
    font=\scriptsize,
    >={Stealth[length=1.6mm]},
    cell/.style={draw, minimum width=5.2mm, minimum height=5mm, inner sep=1pt,
                 align=center, fill=white},
    lab/.style={anchor=east, font=\scriptsize\itshape},
    node distance=2mm,
]

% ---- tiny example DAG (top-left) ----
\node[draw, circle, minimum size=5mm] (A) at (0,0)      {A};
\node[draw, circle, minimum size=5mm] (B) at (1.1,0.7)  {B};
\node[draw, circle, minimum size=5mm] (C) at (1.1,-0.7) {C};
\node[draw, circle, minimum size=5mm] (D) at (2.2,0)    {D};
\draw[->] (A) -- (B); \draw[->] (A) -- (C);
\draw[->] (B) -- (D); \draw[->] (C) -- (D);
\node[font=\scriptsize\itshape] at (1.1,-1.6) {DAG (indices A{=}0,B{=}1,C{=}2,D{=}3)};

% ---- name -> index map ----
\node[lab] (m0) at (5.2, 1.0) {nameToIndex:};
\node[cell] (mA) at (6.1,1.0) {A:0};
\node[cell, right=of mA] (mB) {B:1};
\node[cell, right=of mB] (mC) {C:2};
\node[cell, right=of mC] (mD) {D:3};

% ---- forward CSR ----
\node[lab] (f0) at (5.2, 0.3) {forwardAdj:};
\node[cell] (f1) at (6.1,0.3) {1,2};
\node[cell, right=of f1] (f2) {3};
\node[cell, right=of f2] (f3) {3};
\node[cell, right=of f3] (f4) {--};

% ---- reverse CSR ----
\node[lab] (r0) at (5.2, -0.4) {reverseAdj:};
\node[cell] (r1) at (6.1,-0.4) {--};
\node[cell, right=of r1] (r2) {0};
\node[cell, right=of r2] (r3) {0};
\node[cell, right=of r3] (r4) {1,2};

% ---- topo order / layer ----
\node[lab] (t0) at (5.2, -1.1) {topoOrder:};
\node[cell] (t1) at (6.1,-1.1) {0};
\node[cell, right=of t1] (t2) {1};
\node[cell, right=of t2] (t3) {2};
\node[cell, right=of t3] (t4) {3};
\node[lab] (l0) at (5.2, -1.7) {topoLayer:};
\node[cell] (l1) at (6.1,-1.7) {0};
\node[cell, right=of l1] (l2) {1};
\node[cell, right=of l2] (l3) {1};
\node[cell, right=of l3] (l4) {2};

% ---- ancestor bitsets (rows = variable, columns = bit for A B C D) ----
\node[font=\scriptsize\itshape] at (1.4,-2.6) {ancestorBits[i] (bit set = ancestor)};
\node[lab] at (-0.3,-3.2) {A}; \node[lab] at (-0.3,-3.7) {B};
\node[lab] at (-0.3,-4.2) {C}; \node[lab] at (-0.3,-4.7) {D};
% rows explicitly (A,B,C,D) x (A,B,C,D)
\foreach \row/\y/\vals in {A/-3.2/{0/0,1/0,2/0,3/0}, B/-3.7/{0/1,1/0,2/0,3/0},
                           C/-4.2/{0/1,1/0,2/0,3/0}, D/-4.7/{0/1,1/1,2/1,3/0}} {
  \foreach \c/\v in \vals {
    \node[cell, minimum width=4.5mm] at ({0.1+\c*0.5},\y) {\v};
  }
}

% ---- descendant bitsets ----
\node[font=\scriptsize\itshape] at (5.2,-2.6) {descendantBits[i] (bit set = descendant)};
\node[lab] at (3.6,-3.2) {A}; \node[lab] at (3.6,-3.7) {B};
\node[lab] at (3.6,-4.2) {C}; \node[lab] at (3.6,-4.7) {D};
\foreach \row/\y/\vals in {A/-3.2/{0/0,1/1,2/1,3/1}, B/-3.7/{0/0,1/0,2/0,3/1},
                           C/-4.2/{0/0,1/0,2/0,3/1}, D/-4.7/{0/0,1/0,2/0,3/0}} {
  \foreach \c/\v in \vals {
    \node[cell, minimum width=4.5mm] at ({4.0+\c*0.5},\y) {\v};
  }
}

\end{tikzpicture}
}
  \caption{Compiled index layout for a small fragment. Variables map to dense
  indices; forward/reverse CSR arrays store adjacency; a topological order and
  per-node layer drive pruning; and per-node ancestor/descendant bitsets
  ($\lceil n/64\rceil$ words each) are filled in reverse-/forward-topological
  order. Structural queries become word-parallel bit operations over the
  bitsets.}
  \label{fig:index}
\end{figure}

\subsection{Reachability and path-role queries}\label{sec:index:reach}

With the closures in hand the reachability operators are immediate:
$\mathrm{an}(X)$ and $\mathrm{de}(X)$ are a single bitset read, and
$\textsf{hasDirectedPath}(X,Y)$ is one bit test---is $Y$ set in
$\textsf{descendantBits}[X]$---in $O(1)$. The path-role operators reduce to
bitset algebra over these closures. Confounders of $T$ and $Y$ are the common
ancestors $\textsf{ancestorBits}[T]\cap\textsf{ancestorBits}[Y]$; mediators of
$T\!\to\!Y$ are $\textsf{descendantBits}[T]\cap\textsf{ancestorBits}[Y]$; and
colliders are the nodes whose parents include one reachable from $X$ and one from
$Y$ while the node itself ancestors neither---again a handful of intersections
per candidate. Each is $O(V/64)$ word operations rather than the baseline's
repeated DFS.

\subsection{Conditioning-aware d-separation}\label{sec:index:dsep}

D-separation is \emph{not} reachability: adding a variable to the conditioning set
$Z$ can block one path and, if that variable is a collider or an ancestor of one,
unblock another. The test must therefore follow edge orientation and the role of
each node under $Z$, which we do with the Bayes-Ball procedure. A search
propagates ``balls'' from $X$, tracking for each visited node whether it was
entered from a parent or from a child; at each node the admissible continuations
depend on whether the node is in $Z$: an observed non-collider blocks a ball
arriving along the chain/fork, while a collider (or a node with an observed
descendant) passes a ball arriving from two children. $X$ and $Y$ are
d-separated given $Z$ exactly when no ball reaches $Y$. The procedure is
conditioning-aware by construction and, with conditioning and frontier sets
represented as bitsets, its per-step membership tests are $O(1)$ and the whole
test $O(V/64)$ in the frontier. A single flag---cutting the outgoing edges of the
treatment---turns the same routine into the backdoor test used next.

\subsection{Adjustment-set enumeration}\label{sec:index:adj}

Identification asks whether the effect of $T$ on $Y$ is estimable by covariate
adjustment and, if so, for a valid set. A set $Z$ satisfies the \emph{backdoor
criterion} when it contains no descendant of $T$ and blocks every backdoor path
from $T$ to $Y$; the second condition is exactly the d-separation test of
\S\ref{sec:index:dsep} with $T$'s outgoing edges cut. Enumeration first restricts
the candidate pool to \emph{eligible} variables---those that are not $T$, not $Y$,
and (via $\textsf{descendantBits}[T]$) not descendants of $T$, so the
no-descendant clause holds for free. It then searches candidate sets in
increasing size and returns the first non-empty level of blocking sets, yielding
\emph{minimal} adjustment sets; the empty set is tested first, so a query with no
open backdoor path returns immediately. This is the most traversal-intensive
operator in the catalog---it issues one backdoor (d-separation) test per
candidate set---which is precisely why casting the inner test as an $O(V/64)$
bitset operation, rather than a graph traversal, is what makes identification
scale (R2).

\begin{table}[t]\centering
\caption{Per-operator cost: map-of-sets $+$ DFS baseline vs.\ the bitset index.
$k$ is the number of candidate sets an identification query tests.}
\label{tab:complexity}
\begin{tabular}{@{}lccc@{}}
\toprule
Operator & Baseline query & Index query & Index space \\
\midrule
ancestors / descendants & $O(V{+}E)$ & $O(V/64)$ & $O(V^2/64)$ \\
hasDirectedPath         & $O(V{+}E)$ & $O(1)$     & --- \\
mediators/conf./coll.   & $O(V{+}E)$ & $O(V/64)$  & --- \\
d-separation            & $O(V{+}E)$ & $O(V/64)$  & --- \\
adjustment sets         & $O(k(V{+}E))$ & $O(k\,V/64)$ & --- \\
\bottomrule
\end{tabular}
\end{table}
